% =====================================================================
%  chapters/minggu-12.tex
%  Minggu 12 -- Pengenalan Tailwind CSS
% =====================================================================

\chapter{Pengenalan Tailwind CSS}
\label{chap:mjs-12}

% ---------------------------------------------------------------------
% 1. CAPAIAN PEMBELAJARAN
% ---------------------------------------------------------------------
\begin{capaian}
Setelah menyelesaikan bab ini, mahasiswa diharapkan mampu:
\begin{itemize}
  \item menjelaskan perbedaan antara CSS tradisional dan pendekatan \emph{utility-first};
  \item menyiapkan Tailwind CSS menggunakan \emph{CDN} di proyek HTML;
  \item menggunakan \emph{utility class} dasar untuk tipografi, \emph{spacing}, warna, dan \emph{border};
  \item membuat layout \emph{responsive} menggunakan \emph{Flexbox} dan \emph{Grid} Tailwind.
\end{itemize}
\end{capaian}

% ---------------------------------------------------------------------
% 2. TUJUAN PRAKTIK
% ---------------------------------------------------------------------
\begin{tujuanpraktik}
Pada sesi praktik ini, mahasiswa akan:
\begin{itemize}
  \item mengonversi halaman CSS manual menjadi Tailwind CSS;
  \item bereksperimen dengan \emph{utility class} di \emph{playground};
  \item membangun layout \emph{responsive} dengan \emph{Flexbox} dan \emph{Grid}.
\end{itemize}
\end{tujuanpraktik}

\begin{catatan}
Tailwind CSS \textbf{bukan bahasa pemrograman}---ini adalah kumpulan \textbf{class CSS pendek} yang sudah ditulis untuk Anda. Alih-alih menulis \texttt{font-size: 18px; font-weight: bold;}, Anda cukup menulis \texttt{class="text-lg font-bold"}.
\end{catatan}

% =====================================================================
\section{Konsep: Apa Itu Tailwind CSS?}
\label{sec:12-konsep}
% =====================================================================

\subsection{CSS Tradisional vs Tailwind CSS}

\textbf{CSS Tradisional:}

\begin{lstlisting}[language=CSS, caption={Penulisan CSS tradisional}, label={lst:12-css-tradisional}]
/* Di file terpisah (.css) */
.card {
    background-color: white;
    border-radius: 8px;
    padding: 16px;
    box-shadow: 0 2px 4px rgba(0,0,0,0.1);
}

.card-title {
    font-size: 1.25rem;
    font-weight: bold;
    color: #333;
    margin-bottom: 8px;
}
\end{lstlisting}

\textbf{Dengan Tailwind:}

\begin{lstlisting}[language=HTML, caption={Penulisan dengan Tailwind CSS}, label={lst:12-tailwind}]
<!-- Tidak perlu file CSS terpisah -->
<div class="bg-white rounded-lg p-4 shadow-md">
    <h2 class="text-xl font-bold text-gray-800 mb-2">
        Judul Card
    </h2>
</div>
\end{lstlisting}

\subsection{Mengapa Tailwind?}

\begin{table}[htbp]
\centering
\caption{Kelebihan Tailwind CSS}
\label{tab:12-kelebihan}
\begin{tabular}{ll}
\toprule
\textbf{Kelebihan} & \textbf{Penjelasan} \\
\midrule
Cepat & Tidak perlu bolak-balik ke file CSS \\
Konsisten & \emph{Design system} built-in (spacing, warna, ukuran terstandarisasi) \\
Ringan (\emph{production}) & Hanya class yang dipakai yang masuk ke CSS final \\
Populer & Banyak digunakan di industri \\
\bottomrule
\end{tabular}
\end{table}

\begin{catatan}
CSS tradisional tetap berguna untuk memahami dasar \emph{styling} dan untuk \emph{selector} kompleks yang sulit dijangkau \emph{utility class}. Tailwind \textbf{bukan pengganti} CSS---ini \textbf{pendekatan berbeda} untuk menulis CSS.
\end{catatan}

% =====================================================================
\section{Setup Tailwind di Proyek}
\label{sec:12-setup}
% =====================================================================

Cara termudah adalah menggunakan \textbf{CDN}---tanpa instalasi, cocok untuk tujuan belajar.

\begin{langkah}
  \item Buat file baru \texttt{minggu12.html}.
  \item Buka dengan Live Server. Jika halaman berlatar abu-abu dengan judul besar, Tailwind sudah berhasil dimuat!
\end{langkah}

\begin{lstlisting}[language=HTML, caption={Setup Tailwind CSS via CDN}, label={lst:12-setup}]
<!DOCTYPE html>
<html lang="id">
<head>
    <meta charset="UTF-8">
    <meta name="viewport"
          content="width=device-width, initial-scale=1.0">
    <title>Minggu 12 -- Tailwind CSS</title>
    <!-- Tailwind CSS via CDN -->
    <script src="https://cdn.tailwindcss.com"></script>
</head>
<body class="bg-gray-100 text-gray-800">
    <h1 class="text-3xl font-bold text-center py-10">
        Halo, Tailwind CSS!
    </h1>
</body>
</html>
\end{lstlisting}

\begin{tangkapanlayar}
  {assets/images/minggu-12/tailwind-setup.png}
  {Halaman dengan latar abu-abu terang dan teks judul besar yang sudah ter-\emph{style} oleh Tailwind.}
  {Buka \texttt{minggu12.html} di peramban. Jika judul besar muncul dengan latar abu-abu, setup berhasil.}
\end{tangkapanlayar}

\begin{tip}
Verifikasi dengan DevTools: buka tab Elements $\rightarrow$ \emph{inspect} pada elemen \texttt{<h1>}. Anda akan melihat \emph{class} Tailwind dikonversi menjadi CSS oleh JavaScript CDN.
\end{tip}

\begin{catatan}
\textbf{Catatan CDN:} CDN ini cocok untuk pengembangan dan belajar. Untuk proyek \emph{production}, Tailwind biasanya diinstal via \emph{npm} dan di-\emph{build}.
\end{catatan}

% =====================================================================
\section{Eksplorasi Utility Class Dasar}
\label{sec:12-utility}
% =====================================================================

\subsection{Layout \& Spacing}

\begin{lstlisting}[language=HTML, caption={Utility class untuk spacing}, label={lst:12-spacing}]
<!-- Padding (p = semua sisi, px = horizontal, py = vertical) -->
<div class="p-4">Padding 16px semua sisi</div>
<div class="px-6 py-2">Padding horizontal 24px, vertical 8px</div>

<!-- Margin (m = semua sisi, mx = horizontal, my = vertical) -->
<div class="m-4">Margin 16px semua sisi</div>
<div class="mx-auto">Margin horizontal auto (untuk center) -->

<!-- Skala spacing: angka x 4 = pixel -->
<!-- p-1 = 4px, p-2 = 8px, p-4 = 16px, p-8 = 32px -->
\end{lstlisting}

\subsection{Tipografi}

\begin{lstlisting}[language=HTML, caption={Utility class untuk tipografi}, label={lst:12-tipografi}]
<h1 class="text-4xl">Judul Besar (36px)</h1>
<h2 class="text-2xl">Judul Sedang (24px)</h2>
<p class="text-base">Teks normal (16px)</p>
<p class="text-sm">Teks kecil (14px)</p>

<p class="font-bold">Tebal</p>
<p class="font-semibold">Agak tebal</p>
<p class="font-light">Tipis</p>

<p class="text-red-500">Merah</p>
<p class="text-blue-600">Biru</p>
<p class="text-center">Tengah</p>
\end{lstlisting}

\subsection{Warna \& Background}

\begin{lstlisting}[language=HTML, caption={Utility class untuk warna}, label={lst:12-warna}]
<div class="bg-white">Latar putih</div>
<div class="bg-blue-500">Latar biru</div>
<div class="bg-gray-200">Latar abu-abu terang</div>

<!-- Skala warna Tailwind: 50 (paling terang) sampai 950 (paling gelap) -->
<!-- blue-100 sangat terang, blue-900 sangat gelap -->
\end{lstlisting}

\subsection{Border \& Rounded}

\begin{lstlisting}[language=HTML, caption={Utility class untuk border}, label={lst:12-border}]
<div class="border border-gray-300">Batas abu-abu tipis</div>
<div class="border-2 border-red-500">Batas merah tebal 2px</div>
<div class="rounded-lg">Sudut membulat besar</div>
<div class="rounded-full">Lingkaran penuh</div>
\end{lstlisting}

\subsection{Responsive (Mobile-First)}

\begin{lstlisting}[language=HTML, caption={Kelas responsive Tailwind}, label={lst:12-responsive}]
<!-- Dasar = mobile -->
<!-- sm: >= 640px, md: >= 768px, lg: >= 1024px -->
<div class="text-sm md:text-base lg:text-lg">
    Teks kecil di HP, sedang di tablet, besar di desktop
</div>

<div class="flex flex-col md:flex-row">
    <!-- Vertikal di HP, horizontal di tablet ke atas -->
    <div class="w-full md:w-1/2">Kolom 1</div>
    <div class="w-full md:w-1/2">Kolom 2</div>
</div>
\end{lstlisting}

\subsection{Coba Sendiri: Playground}

\begin{langkah}
  \item Tambahkan kode berikut ke \texttt{minggu12.html} (ganti isi \texttt{<body>}).
  \item \emph{Refresh} dan lihat hasilnya.
\end{langkah}

\begin{lstlisting}[language=HTML, caption={Tailwind \emph{playground}---header, kartu, dan label warna}, label={lst:12-playground}]
<body class="bg-gray-100 text-gray-800 p-6">
    <header class="bg-blue-600 text-white p-4 rounded-lg mb-6">
        <h1 class="text-2xl font-bold">Tailwind Playground</h1>
        <p class="text-blue-100">
            Mencoba utility class dasar
        </p>
    </header>

    <div class="bg-white rounded-lg shadow-md p-6 mb-4">
        <h2 class="text-xl font-semibold text-gray-800 mb-2">
            Kartu Contoh
        </h2>
        <p class="text-gray-600 mb-4">
            Ini adalah kartu dengan padding, shadow,
            dan rounded corners -- semua dari class Tailwind.
        </p>
        <button class="bg-blue-500 hover:bg-blue-700
                       text-white font-bold py-2 px-4 rounded">
            Klik Saya
        </button>
    </div>

    <div class="bg-white rounded-lg shadow-md p-6">
        <h2 class="text-xl font-semibold text-gray-800 mb-2">
            Daftar Warna
        </h2>
        <div class="flex flex-wrap gap-2">
            <span class="bg-red-100 text-red-800
                         px-3 py-1 rounded-full text-sm">
                Red
            </span>
            <span class="bg-blue-100 text-blue-800
                         px-3 py-1 rounded-full text-sm">
                Blue
            </span>
            <span class="bg-green-100 text-green-800
                         px-3 py-1 rounded-full text-sm">
                Green
            </span>
            <span class="bg-yellow-100 text-yellow-800
                         px-3 py-1 rounded-full text-sm">
                Yellow
            </span>
        </div>
    </div>
</body>
\end{lstlisting}

\begin{tangkapanlayar}
  {assets/images/minggu-12/tailwind-playground.png}
  {Halaman dengan kartu, tombol, dan label warna yang sudah di-\emph{style} oleh Tailwind.}
  {Buka \texttt{minggu12.html} di peramban dan bandingkan tampilan Anda.}
\end{tangkapanlayar}

% =====================================================================
\section{Konversi: CSS Manual ke Tailwind}
\label{sec:12-konversi}
% =====================================================================

\begin{langkah}
  \item Buat file \texttt{kartu-profil-css.html} menggunakan CSS manual berikut.
  \item Buat file \texttt{kartu-profil-tailwind.html} menggunakan versi Tailwind.
  \item Bandingkan keduanya di tab browser berbeda.
\end{langkah}

\textbf{Versi CSS Manual:}

\begin{lstlisting}[language=HTML, caption={Kartu profil dengan CSS manual}, label={lst:12-kartu-css}]
<!DOCTYPE html>
<html lang="id">
<head>
    <meta charset="UTF-8">
    <meta name="viewport"
          content="width=device-width, initial-scale=1.0">
    <title>Kartu Profil -- CSS Manual</title>
    <style>
        body {
            font-family: Arial, sans-serif;
            background-color: #f0f2f5;
            display: flex; justify-content: center;
            align-items: center; min-height: 100vh;
            margin: 0;
        }
        .card {
            background-color: white;
            border-radius: 12px; padding: 32px;
            box-shadow: 0 4px 6px rgba(0,0,0,0.1);
            text-align: center; max-width: 320px;
        }
        .avatar {
            width: 100px; height: 100px;
            border-radius: 50%;
            background-color: #4A90D9;
            margin: 0 auto 16px;
            display: flex; align-items: center;
            justify-content: center;
            font-size: 40px; color: white;
        }
        .name {
            font-size: 24px; font-weight: bold;
            color: #333; margin-bottom: 4px;
        }
        .role {
            font-size: 14px; color: #888;
            margin-bottom: 16px;
        }
        .bio {
            font-size: 14px; color: #666;
            line-height: 1.6; margin-bottom: 20px;
        }
        .btn {
            background-color: #4A90D9; color: white;
            border: none; padding: 10px 24px;
            border-radius: 8px; font-size: 14px;
            cursor: pointer;
        }
        .btn:hover { background-color: #357ABD; }
    </style>
</head>
<body>
    <div class="card">
        <div class="avatar">B</div>
        <div class="name">Budi Santoso</div>
        <div class="role">Mahasiswa D3 Informatika</div>
        <div class="bio">
            Siswa semester 1 yang tertarik dengan
            web design dan UI/UX.
        </div>
        <button class="btn">Hubungi</button>
    </div>
</body>
</html>
\end{lstlisting}

\textbf{Versi Tailwind CSS:}

\begin{lstlisting}[language=HTML, caption={Kartu profil dengan Tailwind CSS}, label={lst:12-kartu-tailwind}]
<!DOCTYPE html>
<html lang="id">
<head>
    <meta charset="UTF-8">
    <meta name="viewport"
          content="width=device-width, initial-scale=1.0">
    <title>Kartu Profil -- Tailwind CSS</title>
    <script src="https://cdn.tailwindcss.com"></script>
</head>
<body class="bg-gray-100 flex justify-center
             items-center min-h-screen">
    <div class="bg-white rounded-xl p-8 shadow-lg
                text-center max-w-xs">
        <div class="w-24 h-24 rounded-full bg-blue-500
                    mx-auto mb-4 flex items-center
                    justify-center text-4xl text-white
                    font-bold">
            B
        </div>
        <div class="text-2xl font-bold text-gray-800
                    mb-1">
            Budi Santoso
        </div>
        <div class="text-sm text-gray-400 mb-4">
            Mahasiswa D3 Informatika
        </div>
        <div class="text-sm text-gray-500
                    leading-relaxed mb-5">
            Siswa semester 1 yang tertarik dengan
            web design dan UI/UX.
        </div>
        <button class="bg-blue-500 hover:bg-blue-700
                       text-white font-semibold py-2 px-6
                       rounded-lg transition duration-200">
            Hubungi
        </button>
    </div>
</body>
</html>
\end{lstlisting}

\begin{tangkapanlayar}
  {assets/images/minggu-12/kartu-perbandingan.png}
  {Kartu profil versi Tailwind---tampilan seharusnya sama atau hampir identik dengan versi CSS manual.}
  {Buka kedua file di tab browser berbeda dan bandingkan tampilannya.}
\end{tangkapanlayar}

\subsection{Pemetaan CSS ke Tailwind}

\begin{table}[htbp]
\centering
\caption{Pemetaan properti CSS ke class Tailwind}
\label{tab:12-pemetaan}
\begin{tabular}{ll}
\toprule
\textbf{CSS Manual} & \textbf{Tailwind Equivalent} \\
\midrule
\texttt{background-color: \#f0f2f5;} & \texttt{bg-gray-100} \\
\texttt{display: flex;} & \texttt{flex} \\
\texttt{justify-content: center;} & \texttt{justify-center} \\
\texttt{align-items: center;} & \texttt{items-center} \\
\texttt{min-height: 100vh;} & \texttt{min-h-screen} \\
\texttt{border-radius: 12px;} & \texttt{rounded-xl} \\
\texttt{padding: 32px;} & \texttt{p-8} \\
\texttt{box-shadow: ...;} & \texttt{shadow-lg} \\
\texttt{text-align: center;} & \texttt{text-center} \\
\texttt{max-width: 320px;} & \texttt{max-w-xs} \\
\texttt{width: 100px;} & \texttt{w-24} \\
\texttt{font-size: 24px; font-weight: bold;} & \texttt{text-2xl font-bold} \\
\texttt{color: \#333;} & \texttt{text-gray-800} \\
\texttt{font-size: 14px;} & \texttt{text-sm} \\
\texttt{line-height: 1.6;} & \texttt{leading-relaxed} \\
\texttt{:hover \{...\}} & \texttt{hover:bg-blue-700} \\
\bottomrule
\end{tabular}
\end{table}

% =====================================================================
\section{Layout dengan Tailwind: Flexbox \& Grid}
\label{sec:12-layout}
% =====================================================================

\subsection{Flexbox dengan Tailwind}

\begin{lstlisting}[language=HTML, caption={Flexbox: horizontal dan responsive}, label={lst:12-flexbox}]
<!-- Flex Horizontal (baris) -->
<h3 class="text-lg font-semibold mt-8 mb-3">
    Flex Horizontal (baris)
</h3>
<div class="flex gap-4 bg-white p-4 rounded-lg shadow">
    <div class="bg-blue-500 text-white p-4 rounded-lg
                flex-1 text-center">1</div>
    <div class="bg-green-500 text-white p-4 rounded-lg
                flex-1 text-center">2</div>
    <div class="bg-red-500 text-white p-4 rounded-lg
                flex-1 text-center">3</div>
</div>

<!-- Flex Responsive: vertikal di HP, horizontal di desktop -->
<h3 class="text-lg font-semibold mt-8 mb-3">
    Flex Responsive
</h3>
<div class="flex flex-col md:flex-row gap-4
            bg-white p-4 rounded-lg shadow">
    <div class="bg-purple-500 text-white p-6 rounded-lg
                md:flex-1 text-center">
        Kolom 1
    </div>
    <div class="bg-yellow-500 text-gray-800 p-6 rounded-lg
                md:flex-1 text-center">
        Kolom 2
    </div>
    <div class="bg-teal-500 text-white p-6 rounded-lg
                md:flex-1 text-center">
        Kolom 3
    </div>
</div>
\end{lstlisting}

\subsection{Grid dengan Tailwind}

\begin{lstlisting}[language=HTML, caption={Grid: 3 kolom responsive}, label={lst:12-grid}]
<!-- Grid: 3 kolom responsive -->
<div class="grid grid-cols-1 sm:grid-cols-2 lg:grid-cols-3
            gap-4 bg-white p-4 rounded-lg shadow">
    <div class="bg-indigo-500 text-white p-6 rounded-lg
                text-center">Card 1</div>
    <div class="bg-indigo-500 text-white p-6 rounded-lg
                text-center">Card 2</div>
    <div class="bg-indigo-500 text-white p-6 rounded-lg
                text-center">Card 3</div>
    <div class="bg-indigo-500 text-white p-6 rounded-lg
                text-center">Card 4</div>
    <div class="bg-indigo-500 text-white p-6 rounded-lg
                text-center">Card 5</div>
    <div class="bg-indigo-500 text-white p-6 rounded-lg
                text-center">Card 6</div>
</div>

<!-- Penjelasan grid responsive:
     grid-cols-1     = 1 kolom di HP
     sm:grid-cols-2  = 2 kolom di tablet (>= 640px)
     lg:grid-cols-3  = 3 kolom di desktop (>= 1024px) -->
\end{lstlisting}

\begin{langkah}
  \item \emph{Refresh} halaman dan coba \emph{resize} browser untuk melihat layout \emph{responsive}.
\end{langkah}

\begin{tangkapanlayar}
  {assets/images/minggu-12/tailwind-flexbox.png}
  {Area \emph{flexbox}---tiga kotak berwarna dalam satu baris.}
  {Buka \texttt{minggu12.html} di peramban dan gulir ke bagian Flexbox.}
\end{tangkapanlayar}

\begin{tangkapanlayar}
  {assets/images/minggu-12/tailwind-grid.png}
  {Area \emph{grid}---kotak-kotak \emph{card} dalam format 3 kolom.}
  {Gulir ke bagian Grid dan perhatikan tampilan kolom.}
\end{tangkapanlayar}

% =====================================================================
\section{Latihan}
\label{sec:12-latihan}
% =====================================================================

\begin{latihan}[Konversi Halaman]
Ambil salah satu halaman HTML yang sudah Anda buat di minggu 5--7 (halaman statis dengan CSS). Buat versi baru menggunakan Tailwind CSS.

\begin{enumerate}
  \item Salin file HTML lama ke file baru.
  \item Hapus semua isi \texttt{<style>}.
  \item Tambahkan CDN Tailwind di \texttt{<head>}.
  \item Ganti semua \emph{selector} CSS dengan \emph{class} Tailwind yang sesuai.
\end{enumerate}

\textbf{Output yang diharapkan:}
\begin{itemize}
  \item Halaman terlihat sama atau hampir sama dengan versi CSS manual.
  \item Semua \emph{style} menggunakan \emph{class} Tailwind, tidak ada \texttt{<style>} tersisa.
\end{itemize}

\begin{tangkapanlayar}
  {assets/images/minggu-12/latihan1-konversi.png}
  {Halaman hasil konversi---tampilan sebelah kiri (CSS manual) vs kanan (Tailwind).}
  {Bandingkan halaman asli dengan hasil konversi Tailwind Anda.}
\end{tangkapanlayar}
\end{latihan}

\begin{latihan}[Card Layout Responsive]
Buat file \texttt{latihan12-2.html} yang menampilkan \textbf{3 kartu produk} dalam layout \emph{responsive}:
\begin{enumerate}
  \item Setiap kartu memiliki: gambar \emph{placeholder} (\texttt{bg-gray-300} dengan teks ``Gambar Produk''), nama produk, harga, dan tombol ``Beli''.
  \item Di \emph{mobile}: kartu tampil secara vertikal (1 kolom).
  \item Di tablet ($\geq$ 640\,px): 2 kolom.
  \item Di desktop ($\geq$ 1024\,px): 3 kolom.
\end{enumerate}

Gunakan hanya Tailwind CSS---tidak ada CSS manual.

\begin{tangkapanlayar}
  {assets/images/minggu-12/latihan2-card.png}
  {Kartu produk \emph{responsive}---coba di berbagai ukuran browser.}
  {Buka \texttt{latihan12-2.html} di peramban dan \emph{resize} untuk melihat perubahan kolom.}
\end{tangkapanlayar}
\end{latihan}

\begin{latihan}[Navbar Responsive --- Opsional]
Buat file \texttt{latihan12-3.html} dengan \emph{navbar} yang:
\begin{enumerate}
  \item Horizontal di desktop (logo + navigasi berdampingan).
  \item Di \emph{mobile}: logo di kiri, tombol \emph{hamburger} di kanan, navigasi muncul saat tombol diklik (gunakan JavaScript dari Minggu 11).
  \item Semua \emph{styling} dengan Tailwind CSS.
\end{enumerate}

\begin{tangkapanlayar}
  {assets/images/minggu-12/latihan3-navbar.png}
  {\emph{Navbar} dalam tampilan desktop dan \emph{mobile}.}
  {Buka \texttt{latihan12-3.html} di peramban dan \emph{resize} untuk melihat perubahan layout.}
\end{tangkapanlayar}
\end{latihan}

% =====================================================================
\section{Troubleshooting}
\label{sec:12-troubleshooting}
% =====================================================================

\begin{peringatan}
\textbf{Style Tailwind tidak muncul} --- pastikan \texttt{<script src="https://cdn.tailwindcss.com">} ada di \texttt{<head>}, periksa koneksi internet, periksa ejaan nama \emph{class} (Tailwind sangat spesifik: \texttt{bg-blue-500} $\neq$ \texttt{bg-blue-5000}), dan coba \emph{hard refresh} (\texttt{Ctrl+Shift+R}).
\end{peringatan}

\begin{tip}
\textbf{Tidak tahu nama class yang tepat} --- gunakan \href{https://tailwindcss.com/docs}{Tailwind Docs} sebagai referensi. Cari kategori (mis.\ ``padding'', ``background color'') lalu \emph{copy-paste} \emph{class}. Anda tidak perlu menghafal semua \emph{class}!
\end{tip}

\begin{catatan}
\textbf{Konflik dengan CSS manual} --- jangan campur Tailwind dengan \emph{style} manual yang konflik. Gunakan salah satu: Tailwind SAJA atau CSS manual SAJA.
\end{catatan}

\begin{catatan}
\textbf{CDN lambat atau tidak termuat} --- gunakan versi tertentu (\texttt{https://cdn.tailwindcss.com/3.4.1}). Jika tidak ada internet, kembali ke CSS manual untuk proyek ini.
\end{catatan}

\begin{catatan}
\textbf{Kelas responsive tidak bekerja} --- Tailwind menggunakan \textbf{mobile-first}. Urutan yang benar: \emph{default} (mobile) $\rightarrow$ \texttt{sm} $\rightarrow$ \texttt{md} $\rightarrow$ \texttt{lg} $\rightarrow$ \texttt{xl}.
\end{catatan}
